\documentclass[dvipsnames,hyperref={pdfpagelabels=false}]{beamer}
% beamer already loads hyperref; do NOT load \usepackage{hyperref} again.

\usepackage{ragged2e} % 文字向兩邊對齊
\renewcommand{\raggedright}{\leftskip=0pt \rightskip=0pt plus 0cm}

\usepackage{listings}
\usepackage{lmodern}
\usepackage{fontspec}
\usepackage{xeCJK}
\setCJKmainfont{Noto Sans CJK TC}
\setCJKsansfont{Noto Sans CJK TC}
\setCJKmonofont{Noto Sans Mono CJK TC}
\usepackage[english]{babel}
\usepackage{amsmath,amssymb,amsfonts,amsthm}

% hyperref settings (since beamer already loads hyperref)
\hypersetup{
  colorlinks=true,
  urlcolor=blue,
  citecolor=blue,
  pdfborder={0 0 0}
}

\usepackage{accsupp}% http://ctan.org/pkg/accsupp
\newcommand{\emptyaccsupp}[1]{\BeginAccSupp{ActualText={}}#1\EndAccSupp{}}

\usepackage{booktabs}
\usepackage{colortbl}
\usepackage{xcolor}
\usepackage{threeparttable}

\usepackage{tikz}
\usetikzlibrary{arrows,decorations.pathmorphing,backgrounds,fit,positioning,shapes.symbols,chains}
\usetikzlibrary{tikzmark}

\newcommand{\indep}{\perp\!\!\!\perp}
\newcommand{\nindep}{\not\!\perp\!\!\!\perp}

\beamertemplatefootpagenumber


\title[]
{Guidelines for Oral Exam \\ {\normalsize 口試說明}}

\author
{Prof. Tzu-Ting Yang  \\ 楊子霆}

\institute{Institute of Economics, Academia Sinica  \\ 中央研究院經濟研究所}

\date{\today}


%\usepackage{beamerthemeshadow}  % data (commented: use default beamer theme)
\setbeamertemplate{footline}{%
	\hfill\usebeamercolor[fg]{page number in head/foot}%
	\usebeamerfont{page number in head/foot}%
	\insertframenumber\,/\,\inserttotalframenumber\hspace{1em}\vspace{2pt}%
}

\beamersetuncovermixins{\opaqueness<1>{25}}{\opaqueness<2->{15}}

\begin{document}

\begin{frame}{}
	\titlepage
\end{frame}


%-------------------------------------------------------------
\begin{frame}{Oral Exam: Overview}{口試概述}

\begin{itemize}
	\item 口試可視為 \textbf{term paper 的深度討論}
	\vspace{0.3cm}

	\item 目的是測試兩件事：
	\vspace{0.2cm}
	\begin{itemize}
		\item[(1)] 你在 term paper 中實際使用的\textbf{實證方法}，你是否真正理解？
		\vspace{0.2cm}
		\item[(2)] 你在 homework 中使用的\textbf{程式指令}（R 或 Stata），你是否理解其運作？
	\end{itemize}
	\vspace{0.3cm}

	\item 口試將有\textbf{至少兩個問題}，分別對應上述兩個面向
	\vspace{0.3cm}

	\item 若回答引發進一步疑問，教授將當場追問 follow-up 問題
\end{itemize}

\end{frame}


%-------------------------------------------------------------
\begin{frame}{Question 1: Empirical Method}{第一題：實證方法}

\begin{block}{測試內容}
	你在 term paper 中使用的\textbf{實證方法}，需能當場：
	\begin{itemize}
		\vspace{0.1cm}
		\item \textbf{寫出數學式}（例如：迴歸方程式、識別假設、估計量公式）
		\vspace{0.1cm}
		\item \textbf{解釋每個符號的意義}，以及整個方法的直覺
		\vspace{0.1cm}
		\item 回答關於該方法的 \textbf{follow-up 問題}
	\end{itemize}
\end{block}

\vspace{0.3cm}

\begin{exampleblock}{準備建議}
	\begin{itemize}
		\item 重讀你 term paper 的 empirical strategy 段落
		\vspace{0.1cm}
		\item 確認你能\textbf{不看投影片}，從頭到尾口頭說明你的方法
		\vspace{0.1cm}
		\item 思考：這個方法的\textbf{關鍵假設}是什麼？若假設不成立會怎樣？
	\end{itemize}
\end{exampleblock}

\end{frame}


%-------------------------------------------------------------
\begin{frame}{Question 1: Examples}{第一題：範例}

\vspace{0.1cm}
\small
以下為可能被問到的問題範例（依你的 term paper 方法而定）：

\vspace{0.3cm}
\begin{tabular}{p{0.30\textwidth} p{0.62\textwidth}}
\toprule
\textbf{方法} & \textbf{可能被問} \\
\midrule
OLS / 迴歸 &
請寫出你的迴歸方程式，並解釋係數 $\hat{\beta}$ 的意義 \\[6pt]
Difference-in-Differences &
請寫出 DiD 的識別假設，平行趨勢假設是什麼意思？ \\[6pt]
Instrumental Variables &
請寫出 2SLS 的第一階段，你的工具變數為什麼有效？ \\[6pt]
Regression Discontinuity &
請說明 RDD 的識別策略，斷點前後的估計量代表什麼？ \\[6pt]
Matching / PSM &
請說明 propensity score 的定義，匹配的直覺是什麼？ \\
\bottomrule
\end{tabular}

\end{frame}


%-------------------------------------------------------------
\begin{frame}{Question 2: Programming Code}{第二題：程式碼}

\begin{block}{測試內容}
	根據你在 homework 使用的程式語言（\textbf{R} 或 \textbf{Stata}），需能解釋：
	\begin{itemize}
		\vspace{0.1cm}
		\item 某段\textbf{程式碼的功能}是什麼？每個指令或選項代表什麼意思？
		\vspace{0.1cm}
		\item 輸出結果（output）如何\textbf{解讀}？
		\vspace{0.1cm}
		\item 若改變某個參數，\textbf{結果會如何變化}？
	\end{itemize}
\end{block}

\vspace{0.3cm}

\begin{exampleblock}{準備建議}
	\begin{itemize}
		\item 重新執行你的 homework 程式碼，確認每一行你都看得懂
		\vspace{0.1cm}
		\item 不需要背指令，但需要理解\textbf{每個指令在做什麼}
		\vspace{0.1cm}
		\item 重點在於：看到一段程式碼，\textbf{能說出它對應到哪個統計概念}
	\end{itemize}
\end{exampleblock}

\end{frame}


%-------------------------------------------------------------
\begin{frame}{Question 2: Examples}{第二題：範例}

\vspace{0.1cm}
\small
以下為可能被問到的程式碼問題範例：

\vspace{0.25cm}
\begin{block}{R 範例}
\texttt{feols(y \textasciitilde{} x1 + x2 | id + year, data = df, cluster = \textasciitilde{}id)}
\begin{itemize}
	\item \texttt{| id + year} 代表什麼？為什麼需要加 fixed effects？
	\item \texttt{cluster = \textasciitilde{}id} 的作用是什麼？不加會怎樣？
\end{itemize}
\end{block}

\vspace{0.2cm}
\begin{block}{Stata 範例}
\texttt{ivregress 2sls y x1 (x2 = z), robust first}
\begin{itemize}
	\item \texttt{(x2 = z)} 的意思是什麼？\texttt{z} 扮演什麼角色？
	\item \texttt{first} 這個選項輸出什麼？你如何判斷工具變數夠強？
\end{itemize}
\end{block}

\end{frame}


%-------------------------------------------------------------
\begin{frame}{Summary}{總結}

\begin{center}
\begin{tikzpicture}[
  box/.style={rectangle, draw, rounded corners, align=center,
              minimum width=3.8cm, minimum height=1.1cm,
              font=\small, fill=blue!8},
  qbox/.style={rectangle, draw, rounded corners, align=center,
               minimum width=5.0cm, minimum height=1.3cm,
               font=\small},
  arrow/.style={->, >=stealth, thick}
]

\node[box] (tp) at (0, 2.2) {Term Paper\\討論 \& 深度理解};

\node[qbox, fill=orange!20] (q1) at (-3.0, 0.4)
  {\textbf{Q1: 實證方法}\\寫出數學式\\解釋意義 + follow-up};

\node[qbox, fill=green!15] (q2) at (3.0, 0.4)
  {\textbf{Q2: 程式碼}\\R 或 Stata\\解釋指令 \& 輸出};

\draw[arrow] (tp) -- (q1);
\draw[arrow] (tp) -- (q2);

\node[font=\scriptsize, align=center] at (-3.0, -0.75)
  {依據你的 term paper 方法};
\node[font=\scriptsize, align=center] at (3.0, -0.75)
  {依據你的 homework 程式語言};

\end{tikzpicture}
\end{center}

\vspace{0.1cm}
\begin{center}
\fbox{\parbox{0.85\textwidth}{\centering
	口試的目的不是為難你，而是確認你真正理解\\
	自己在 term paper 和 homework 中做了什麼。
}}
\end{center}

\end{frame}


\end{document}
